\section{Propiedades específicas de la puerta de Pauli-\gate{X} controlada}
En esta sección se demuestran algunas propiedades específicas que involucran a la puerta de Pauli-\gate{X} controlada. En todas las demostraciones $\ket{a}$ es un ket producto, que se descompone según los valores de los cúbits que intervienen mediante~\eqref{eq:descomposicion-indices}; como las puertas son lineales, basta calcular su acción sobre cada componente.

Para $x,y\in\B$, la definición~\eqref{eq:generic-controlled-gate} da directamente
\begin{equation}
  \label{eq:cnot-base}
  \stategate{X}{h}{k}\stateket{a}{hk}{xy}=\stateket{a}{hk}{x\,(x\oplus y)},
\end{equation}
es decir, la puerta $\stategate{X}{h}{k}$ suma al cúbit $k$ el valor del cúbit $h$.

\begin{proposition}
  Sean $h,k \in \mathbb{N}_n$ dos índices distintos.
  Entonces
  \begin{equation}
    \stategate{X}{h}{k}\stategate{Z}{}{k}\stategate{X}{h}{k} = \stategate{Z}{}{h,k}
  \end{equation}
\end{proposition}
\begin{proof}
  Sea $\ket{a}=a_{h0}\stateket{a}{h}{0}+a_{h1}\stateket{a}{h}{1}$. Sobre la componente $\stateket{a}{h}{0}$ las dos puertas controladas actúan como la identidad. Sobre $\stateket{a}{h}{1}$ actúan como $\gate{X}_k$, porque $\gate{Z}_k$ no modifica el cúbit $h$, y $\gate{X}\gate{Z}\gate{X}=-\gate{Z}$. Por tanto
  \begin{align*}
    \stategate{X}{h}{k}\stategate{Z}{}{k}\stategate{X}{h}{k}\ket{a}
     & =a_{h0}\,\gate{Z}_k\stateket{a}{h}{0}+a_{h1}\,\gate{X}_k\gate{Z}_k\gate{X}_k\stateket{a}{h}{1}=   \\
     & =a_{h0}\,\gate{Z}_k\stateket{a}{h}{0}-a_{h1}\,\gate{Z}_k\stateket{a}{h}{1}=                      \\
     & =\gate{Z}_k\left(a_{h0}\stateket{a}{h}{0}-a_{h1}\stateket{a}{h}{1}\right)=\gate{Z}_k\gate{Z}_h\ket{a}
    =\stategate{Z}{}{h,k}\ket{a},
  \end{align*}
  donde en el penúltimo paso se ha usado la expresión de $\gate{Z}_h$ de la sección de puertas de Pauli.
\end{proof}

Para el siguiente resultado es necesario introducir la puerta SWAP, que actúa sobre dos cúbits intercambiando el estado de ambas posiciones:
\begin{align}
  \label{eq:swap}
  \begin{split}
    (\gate{SWAP})_{h,k}\ket{a} & =a_{h0}a_{k0}\stateket{a}{hk}{00}+a_{h0}a_{k1}\stateket{a}{hk}{10}+ \\
                               & +a_{h1}a_{k0}\stateket{a}{hk}{01}+a_{h1}a_{k1}\stateket{a}{hk}{11}
  \end{split}
\end{align}
\begin{proposition}
  Sean $h,k \in \mathbb{N}_n$ dos índices distintos.
  Entonces
  \begin{equation}
    \stategate{X}{h}{k}\stategate{X}{k}{h}\stategate{X}{h}{k} = (\gate{SWAP})_{h,k}
  \end{equation}
\end{proposition}
\begin{proof}
  Sean $x,y\in\B$. Aplicando~\eqref{eq:cnot-base} tres veces (y su versión con los papeles de $h$ y $k$ intercambiados),
  \[
    \stategate{X}{h}{k}\stategate{X}{k}{h}\stategate{X}{h}{k}\stateket{a}{hk}{xy}
    =\stategate{X}{h}{k}\stategate{X}{k}{h}\stateket{a}{hk}{x\,(x\oplus y)}
    =\stategate{X}{h}{k}\stateket{a}{hk}{y\,(x\oplus y)}
    =\stateket{a}{hk}{yx}.
  \]
  Por~\eqref{eq:descomposicion-indices} y la linealidad,
  \[
    \stategate{X}{h}{k}\stategate{X}{k}{h}\stategate{X}{h}{k}\ket{a}
    =\sum_{x,y\in\B}a_{hx}a_{ky}\stateket{a}{hk}{yx},
  \]
  que es exactamente~\eqref{eq:swap}.
\end{proof}

\begin{proposition}
  Sean $h,k \in \mathbb{N}_n$ dos índices distintos.
  Entonces
  \begin{equation}
    \gate{H}_{h,k}\stategate{X}{k}{h}\gate{H}_{h,k} = \stategate{X}{h}{k}
  \end{equation}
\end{proposition}
\begin{proof}
  Para $x\in\B$ sea $s(x)$ el signo $+$ si $x=0$ y $-$ si $x=1$, de modo que $\gate{H}\ket{x}=\ket{s(x)}=\frac{1}{\sqrt2}\left(\ket{0}+(-1)^x\ket{1}\right)$ y, como $\gate{H}\gate{H}=\gate{I}$, también $\gate{H}\ket{s(x)}=\ket{x}$. Además $\gate{X}\ket{s(x)}=(-1)^x\ket{s(x)}$.

  Sean $x,y\in\B$. En primer lugar,
  \begin{align*}
    \stategate{X}{k}{h}\stateket{a}{hk}{s(x)\,s(y)}
     & =\tfrac{1}{\sqrt2}\left(\stategate{X}{k}{h}\stateket{a}{hk}{s(x)\,0}+(-1)^y\stategate{X}{k}{h}\stateket{a}{hk}{s(x)\,1}\right)= \\
     & =\tfrac{1}{\sqrt2}\left(\stateket{a}{hk}{s(x)\,0}+(-1)^{y}(-1)^{x}\stateket{a}{hk}{s(x)\,1}\right)
    =\stateket{a}{hk}{s(x)\,s(x\oplus y)},
  \end{align*}
  porque sobre la componente en la que el cúbit $k$ vale $1$ se aplica $\gate{X}_h$ a $\ket{s(x)}$. Por tanto
  \[
    \gate{H}_{h,k}\stategate{X}{k}{h}\gate{H}_{h,k}\stateket{a}{hk}{xy}
    =\gate{H}_{h,k}\stategate{X}{k}{h}\stateket{a}{hk}{s(x)\,s(y)}
    =\gate{H}_{h,k}\stateket{a}{hk}{s(x)\,s(x\oplus y)}
    =\stateket{a}{hk}{x\,(x\oplus y)},
  \]
  que por~\eqref{eq:cnot-base} es $\stategate{X}{h}{k}\stateket{a}{hk}{xy}$. Por~\eqref{eq:descomposicion-indices} y la linealidad se obtiene la igualdad sobre $\ket{a}$.
\end{proof}

\begin{proposition}
  Sean $j,h,k,l \in \mathbb{N}_n$ cuatro índices distintos.
  Entonces
  \begin{equation}
    \label{eq:mainresult}
    \stategate{X}{jhk}{l}\gate{X}_h\stategate{X}{jhk}{l} = \stategate{X}{jk}{l}\gate{X}_h
  \end{equation}
\end{proposition}
\begin{proof}
  Sean $x,y,z\in\B$ y consideremos la componente $\stateket{a}{jhk}{xyz}$. Se aplican las tres puertas del lado izquierdo, de derecha a izquierda:
  \begin{itemize}
    \item la primera $\stategate{X}{jhk}{l}$ aplica $\gate{X}_l$ si y solo si $x=y=z=1$;
    \item $\gate{X}_h$ cambia el valor del cúbit $h$ de $y$ a $1\oplus y$;
    \item la segunda $\stategate{X}{jhk}{l}$ aplica $\gate{X}_l$ si y solo si $x=z=1$ y $1\oplus y=1$, es decir, si $x=z=1$ e $y=0$.
  \end{itemize}
  Ninguna de estas puertas modifica los cúbits $j$ y $k$. Si $x=z=1$, exactamente una de las dos puertas controladas aplica $\gate{X}_l$ (la primera si $y=1$ y la segunda si $y=0$); en otro caso no la aplica ninguna. Como $\gate{X}_l$ y $\gate{X}_h$ conmutan por~\eqref{eq:conmutacion-gates},
  \[
    \stategate{X}{jhk}{l}\gate{X}_h\stategate{X}{jhk}{l}\stateket{a}{jhk}{xyz}=
    \begin{cases}
      \gate{X}_l\gate{X}_h\stateket{a}{jhk}{xyz} & \text{si } x=z=1, \\
      \gate{X}_h\stateket{a}{jhk}{xyz}           & \text{en otro caso,}
    \end{cases}
  \]
  que es precisamente $\stategate{X}{jk}{l}\gate{X}_h\stateket{a}{jhk}{xyz}$, ya que $\gate{X}_h$ no cambia los cúbits de control $j$ y $k$. Por~\eqref{eq:descomposicion-indices} y la linealidad se obtiene la igualdad sobre $\ket{a}$.
\end{proof}
